\section{Finite coefficient extension and the actual standard subring}
\label{sec:finite-base-change}

\status{BC1--BC8 construct the source's standard coefficient ring at a
specified finite place and prove the actual \texttt{LPBaseChange} theorem.
The parameters $t,D,\kappa,c_*$ are unchanged. This closes the finite
extension step that HM12's ring-isomorphism argument did not address.
For a literal \texttt{recC} datum the standard subring is automatic;
existence of the full globally compatible family remains separate.}

\paragraph{BC1: what the standard ring actually is.}
Let $p\ge5$ be prime, $f\ge2$ even and $q=p^f$. With $W(k)$ denoting
Witt vectors, set
\[
 W=W(\mathbb F_q),\quad W_2=W(\mathbb F_{q^2}),\quad
 O_s=W_2[\omega]/(\omega^2-p),\quad K_s=\operatorname{Frac}(O_s).
                                                               \tag{BC.1}
\]
The field $Q=\operatorname{Frac}(W)$ embeds in $K_s$. Its Frobenius
$\sigma$ reduces to $x\mapsto x^p$ and satisfies $\sigma^f=1$.
The extension $W_2/W$ is unramified of degree two, whereas adjoining
$\omega$ is ramified of degree two. Thus $O_s$ is free of rank four
over $W$, has uniformizer $\omega$ and residue field $\mathbb F_{q^2}$.
These are the source's \texttt{WQ, OO, KO, jO, iota, wU}.

To check the assertions directly, lift a residue-field basis of
$\mathbb F_{q^2}/\mathbb F_q$ and successively subtract its residue
expansion, dividing the remainder by $p$. Completeness gives a spanning
basis of $W_2$ over $W$; reduction of a relation after removing its
smallest $p$-power proves independence. The polynomial $X^2-p$ is
irreducible since a square in $\operatorname{Frac}(W_2)$ has even
$p$-valuation. For $a+b\omega\ne0$ the candidate orders
$2v_p(a)$ and $2v_p(b)+1$ differ. Factoring the smaller power of
$\omega$ leaves a unit. This proves the DVR and residue assertions;
$\omega^2=p$ also proves completeness.

Every integer is a square in $O_s$. Write $z=p^m u$ with $p\nmid u$.
Every nonzero element of $\mathbb F_q$ is a square in
$\mathbb F_{q^2}$, since $q-1$ divides $(q^2-1)/2$.
Hensel iteration lifts a square root of $\bar u$ to $y\in W_2$
(the derivative $2y$ is a unit), and $(\omega^m y)^2=z$.
The case $z=0$ is immediate. This checks the full standard quadratic
setting, including its finite structural map from $W$.

\paragraph{BC2: the finite extension at the chosen place exists.}
We need two elementary facts, with proofs recorded here because they
supply the actual coefficient ring.

First, if $R$ is a complete DVR with characteristic-zero fraction field
$K$ and $E/K$ is finite, its integral closure $S$ in $E$ is a finite
free complete DVR. Choose an integral primitive element and the full
order $A=R[\theta]$. Separability and integrality of traces give
\[
 A\subset S\subset A^\#,
 \qquad A^\#=\{x\in E:\operatorname{Tr}_{E/K}(xA)\subset R\}.
                                                               \tag{BC.2}
\]
The trace form is nondegenerate, so $A^\#$ is a finite fractional
lattice. Hence $S$ is finite and torsion free, therefore free over $R$.
Clearing polynomial coefficients shows $\operatorname{Frac}(S)=E$.
The ring $S$ is complete in the uniformizer-adic topology. If it had
several maximal ideals, the finite Artinian ring $S/\pi S$ would
have a nontrivial idempotent. Elements of $1+\pi S$ are units by the
convergent geometric series; Newton iteration for $X^2-X$ therefore lifts it
to $S$, contradicting that $S$ is a domain. Thus $S$ is local.
It is normal and one-dimensional, hence a DVR; explicitly, for
$0\ne a\in\mathfrak m$, the equality $\sqrt{(a)}=\mathfrak m$
and finite generation give a power $\mathfrak m^n\subset(a)$. Choose the least $n$ with
$\mathfrak m^n\subset(a)$ and $b\in\mathfrak m^{n-1}\setminus(a)$.
Then $(b/a)\mathfrak m\subset S$. If this lay in $\mathfrak m$,
the determinant trick would make $b/a$ integral. Otherwise
$(b/a)x$ is a unit for some $x\in\mathfrak m$, forcing
$\mathfrak m=(x)$. The two adic topologies agree, proving completeness.

Second, let $L$ be the source number field and $v$ a specified prime
above $p$. Inside $\Omega=\overline{K_s}$, each embedding $\tau:L\to\Omega$
defines
\[
 E_\tau=K_s\tau(L),\qquad
 S_\tau=\text{integral closure of }O_s\text{ in }E_\tau.  \tag{BC.3}
\]
This extension is finite, and its maximal ideal induces a prime
$P_\tau$ of $L$ above $p$. Suppose none is $v$. The Chinese remainder
theorem gives an integer $a\in L$ that is $0$ modulo $v$ and $1$
modulo every distinct $P_\tau$. All $\tau(a)^{-1}$ are integral over
$O_s$, so their product $N_{L/\mathbb Q}(a)^{-1}$ is integral.
But $p\mid N_{L/\mathbb Q}(a)\ne0$, since the norm is an integer
multiple of $a$ in the ring of integers. This makes $p^{-1}$ integral
over the integrally closed ring $O_s$, a contradiction.
Thus some $\tau$ induces $v$, and
\[
 \tau(x)\in S_\tau\quad\Longleftrightarrow\quad v(x)\ge0.
                                                               \tag{BC.4}
\]
For the reverse implication use a local fraction with denominator
outside $v$; for the converse use $x^{-1}$ in the local maximal ideal.
These are precisely the objects \texttt{KfRec, ORec, embRec, extRec}.

\paragraph{BC3: exact lattice identities and ramification.}
Let $\phi:R\hookrightarrow S$ be a finite inclusion of DVRs, with
fraction fields $K\hookrightarrow K'$. The finite torsion-free
$R$-module $S$ is free. Moreover $1$ is primitive: if $\pi$ were
invertible in $S$, then $\pi^{-1}\in K$ would be integral over $R$.
Choose an $R$-basis $b_1=1,b_2,\ldots,b_d$ of $S$. For a
$K$-space $V$ and any $R$-submodule $M\subset V$,
\[
 SM=\bigoplus_{a=1}^d b_aM\subset K'\otimes_KV.
 \quad\text{Consequently}\quad
 S\Bigl(\bigcap_iM_i\Bigr)=\bigcap_iSM_i,
 \qquad SM\cap V=M.                                     \tag{BC.5}
\]
The intersection may be infinite: finite freeness, not flatness alone,
justifies the coordinate argument. Kernels, invariant submodules defined
by universal coefficient equations, saturation and ranks also commute
with this extension. These statements give faithful descent of lattice
inclusions and equalities.

The dual map is the semilinear map
$\delta(G)(a\otimes x)=a\phi_K(G(x))$. A lattice basis and its
algebraic dual basis prove
\[
 S\delta(M^\vee)=(SM)^\vee.                              \tag{BC.6}
\]
The pairing is $K'$-valued evaluation. No trace-dual or different is
involved. If $\phi(\pi_R)=u\pi_S^e$, then
\[
 \operatorname{ord}_S\phi_K(x)=e\operatorname{ord}_R(x),
 \qquad
 \frac{\operatorname{ord}_S\phi_K(x)}{\operatorname{ord}_S(p)}
 =\frac{\operatorname{ord}_R(x)}{\operatorname{ord}_R(p)}.
                                                               \tag{BC.7}
\]
Lengths multiply by $e$, as follows by applying extension to each
Smith factor $R/(\pi_R^a)$. Normalized lengths and determinant
valuations are unchanged. In particular the image of $\omega$ has
integer valuation $e$ and $p$-normalized valuation $1/2$; it need
not be a uniformizer.

\paragraph{BC4--BC5: identify the actual polynomial lattices.}
CB2 supplies the same monic integral tableau basis for the literal
dual-Weyl module $\nabla_\mu$ over every domain. Hence
$S\otimes_R\nabla_\mu(R)=\nabla_\mu(S)$, including the generic
unipotent-invariance identities. For decreasing integral $\lambda$,
put $\mu_i=\lambda_i-\lambda_k$ and $T(g)f(X)=f(gX)$. The actual
lattice at $\Lambda=gR^k$ is
\[
 H_\lambda(\Lambda)
 =\det(g)^{-\lambda_k}T(g^{-1})\nabla_\mu(R).             \tag{BC.8}
\]
Replacing $g$ by $gh$, $h\in\mathrm{GL}_k(R)$, leaves it unchanged:
$T((gh)^{-1})=T(g^{-1})T(h^{-1})$, the inner factor preserves
$\nabla_\mu(R)$, and $\det(h)^{-\lambda_k}$ is a unit.
Thus the source's intersection over all lattice bases equals
(BC.8). Determinants, inverses and substitution commute with coefficient
maps, so (BC.8) commutes with extension. Products of tableau bases in
disjoint variable blocks give the same assertion for the actual tensor.

Here is the full dictionary. For $s=0,\ldots,f-1$, use
$\lambda_s=\lambda_s^0$ when $s$ is even and
$\lambda_{s,i}=-\lambda^0_{s,k+1-i}$ when $s$ is odd. Set
$A_r=\sum_s\sum_{i\le r}\lambda_{s,i}$ and
\[
\begin{aligned}
 \Lambda_s(L)&=O\,\iota(\sigma^sL),&
 \Lambda_s(e)&=\Lambda_s(e.\mathrm{src})
                   +\omega^{-1}\Lambda_s(p\,e.\mathrm{tgt}),\\
 H(L)&=\bigotimes_s H_{\lambda_s}(\Lambda_s(L)),&
 H(e)&=\bigotimes_s H_{\lambda_s}(\Lambda_s(e)),\\
 R_e&=\omega^{-A_r}H(e),&R'_e&=\omega^{-A_r}H(e)^\vee,
 \qquad b(F,G)=p^{A_r}G(F).
\end{aligned}                                             \tag{BC.9}
\]
An adapted $W$-basis of the edge has target
$g\operatorname{diag}(p^{-1}I_r,I)W^k$ and source $gW^k$;
therefore its midpoint lattice is
$\iota(\sigma^sg)\operatorname{diag}(\omega^{-1}I_r,I)O^k$.
These are genuine full lattices, with no implicit choice of representation.

Let $\nu$ map polynomial coefficients and $\delta$ map duals. Their
common tableau basis identifies the two representation spaces by scalar
extension, with equal dimension. The exact clean lattices are
\[
\begin{aligned}
 N_{\rm src}(L)&=\omega^{-t}H(L)\cap
                           \bigcap_{e.\mathrm{src}=L}R_e,\\
 N_{\rm tgt}(L)&=\omega^{-t}H(L)^\vee\cap
                           \bigcap_{e.\mathrm{tgt}=L}R'_e.
\end{aligned}                                             \tag{BC.10}
\]
Equations (BC.5)--(BC.9) identify their scalar extensions with the
actual target clean lattices. They also preserve every extra
$\omega^{c_*}R_{e'}$ condition. Crucially,
$b'(\nu F,\delta G)=\phi_K(b(F,G))$, even if $A_r<0$.

\paragraph{BC6--BC7: transport the rank certificate without loss.}
The source's \texttt{GoodRankK} certificate is an integral map
$u:O^n\to N_{\rm tgt}(L)$ and an integral matrix $M$ with
\[
 \sum_i M_{ji}b(F_i,u(x))=p^Dx_j\qquad(x\in O^n).
\]
For $y\in S^n$ put $u'(y)=\sum_\ell y_\ell\delta(u(\epsilon_\ell))$
and $M'=\phi(M)$. The previous lattice identities give the correct
codomain, and
\[
 \sum_iM'_{ji}b'(\nu F_i,u'(y))
 =\sum_\ell y_\ell\phi_K\!\left(
         \sum_iM_{ji}b(F_i,u(\epsilon_\ell))\right)=p^Dy_j.
                                                               \tag{BC.11}
\]
Keep the same $W$-edges and send $F_i$ to $\nu F_i$. All clean
and alternative-edge divisibilities hold, and the rank inequality
$n\ge(1-\kappa)\dim V$ is unchanged. So is $4(D+1)\le c_*$.
This proves
\[
 \texttt{LPprime}_{O}(t,D,\kappa,c_*)
 \ \Longrightarrow\
 \texttt{LPprime}_{S}(t,D,\kappa,c_*).                    \tag{BC.12}
\]
If a source weight is not dominant, the defined representation has a
zero tensor factor; the empty list and $n=0$ supply the literal
exceptional branch. The larger coefficient residue field never replaces
$q=\#(W/pW)$ in the edge data. The expanded report also proves descent
for a fixed family of tests by fractional Smith form; no arbitrary newly
chosen family of vectors is claimed to descend.

\paragraph{BC8: match the actual global-family interface.}
The source \texttt{StdSubAt} asks for the standard ring and a finite map
satisfying $\phi_K\iota_s=\iota$, $\phi j_s=j$ and
$\phi(\omega_s)=\omega$. In \texttt{locPackC/recC}, the rings are
literally $WQ,QQ,ORec,KfRec$, and all three maps are the composites
constructed in BC1--BC2. Thus \texttt{ringTied\_recC} uses identity
isomorphisms and the interface follows. For a compatibly isomorphic
presentation, compose the field/ring isomorphisms and their inverses;
finiteness and the three displayed equations persist. The actual
Frobenius on $Q$ stays fixed, using its supplied period and reduction
properties.

This proves \texttt{StdSubAt} pointwise for every well-formed
\texttt{recC}, and eventually for a supplied \texttt{IsRecCEv} family.
The latter is one sufficient identification, not a necessary route:
an actual \texttt{GroupCompatibilityFamily} also supplies the required
ring identifications directly, and BC applies to its ring field.
It does not assume that an arbitrary Hida coefficient DVR already
contains the doubled residue field: that field was deliberately built
into \texttt{recC}. Supplying the remaining global construction and its
simultaneous compatibility is the precise boundary retained here.
